\subsection{Empirical Validation of Residual Suppression}
\label{sec:residual-suppression-ablation}

We test the one-way implication of
Proposition~\ref{prop:residual-gating} across CIFAR-10, CIFAR-100,
Tiny-ImageNet, and SVHN.  For each dataset, we use ten targets and five
independently trained models, yielding 50 target--model pairs.  Within each
pair, we order all adversarial-class candidates by their signed margin $M_i$
and divide them into 30 equal-sized groups.  We then record the residual norm
$\lVert r_i\rVert_2$, the magnitude of the classifier-head interaction
$|H_i|=|(r_i^\top r_t)(h_i^\top h_t)|$, and the magnitude of the backbone
interaction
$|A_i|=|\langle J_i^\top W^\top r_i,
J_t^\top W^\top r_t\rangle|$.  This produces 200 target--model curves and
532,950 candidate evaluations in total.  Each thin gray curve in
Figure~\ref{fig:residual-suppression-multidataset} reports the median within
one target--model pair, while the thick blue curve reports the overall median
at the same margin group.

\begin{figure*}[t]
    \centering
    \includegraphics[width=\textwidth]{residual_suppression_30_groups_multidataset.pdf}
    \caption{
    Empirical validation of residual suppression on four datasets.  Candidates
    are ordered from low to high signed margin separately for every
    target--model pair and partitioned into 30 equal-sized groups.  The columns
    show the candidate residual norm, the absolute classifier-head interaction,
    and the absolute backbone interaction.  Thin gray curves correspond to
    individual target--model pairs, and the thick blue curve is the overall
    median.  All vertical axes use a logarithmic scale.
    }
    \label{fig:residual-suppression-multidataset}
\end{figure*}

The same pattern appears in every dataset.  As the signed margin increases,
the residual norm rapidly decreases, and both gradient-interaction terms
decrease with it.  The backbone column is especially important because it
measures the Jacobian-dependent term directly rather than inferring its
behavior from representation similarity.  These results support the mechanism
behind Proposition~\ref{prop:residual-gating}.  A candidate with a sufficiently
large positive margin has a negligible residual, which suppresses both the
classifier-head and backbone interactions even when its feature or Jacobian
geometry might otherwise be favorable.  The implication remains one-way.
Avoiding this suppressed regime does not guarantee positive full-gradient
alignment or a successful poison, but it removes candidates for which both
local alignment channels are already weak.
